\documentclass[conference,10pt]{IEEEtran}
\usepackage{amsmath,amssymb,amsthm,amsfonts}
\usepackage{booktabs}
\usepackage{multirow}
\usepackage{graphicx}
\usepackage{hyperref}
\usepackage{xcolor}
\usepackage{microtype}

\newtheorem{theorem}{Theorem}
\newtheorem{proposition}{Proposition}
\newtheorem{corollary}{Corollary}

\title{Dual-NVFP4: Surpassing Native BF16 Speed and MXFP8 Accuracy on Blackwell \texttt{tcgen05} via Residual-Microscaled FP4 and SMEM-Reused Tensor Cores}

\author{\IEEEauthorblockN{Rishi Vinodkumar}}

\begin{document}
\maketitle

\begin{abstract}
NVIDIA Blackwell (\texttt{GB200}, \texttt{sm\_100a}) delivers $4\times$ higher theoretical Tensor-Core throughput in block-scaled \texttt{NVFP4} ($1{\times}16$ \texttt{E4M3} micro-scales, $9.65\text{ PF/s}$ peak, $6.56$--$7.23\text{ PF/s}$ live \texttt{CUTLASS}) than in \texttt{BF16} ($2.41\text{ PF/s}$ peak, $1.54$--$1.85\text{ PF/s}$ empirical random \texttt{cuBLAS}). However, single-term \texttt{NVFP4} suffers severe quantization degradation under $50\times$--$200\times$ LLM activation outliers and multi-layer MoE routing drift, and standard Hadamard rotation (\texttt{FWHT}) \emph{worsens} \texttt{NVFP4} on SwiGLU LLMs by destroying fine-grained $1{\times}16$ sparsity (increasing 48-layer \texttt{Qwen3-30B-A3B} KL divergence by $2.15\times$).

We introduce \textbf{\texttt{Dual-NVFP4}}, a rotation-free residual micro-scaled \texttt{FP4} decomposition $A = A_0 + A_1$ where $A_0 = Q_{\text{NVFP4}}(A)$ absorbs intra-block outliers and the residual $A_1 = Q_{\text{NVFP4}}(A - A_0)$ receives its own independent $1{\times}16$ \texttt{E4M3} scale $s_1 \ll s_0$. Dropping the $O(2^{-8})$ cross-term $A_1 B_1$ reduces matrix multiplication to $3$ \texttt{NVFP4} Tensor-Core MMAs ($A_0 B_0 + A_1 B_0 + A_0 B_1$) accumulated directly into Blackwell Tensor Memory (\texttt{TMEM}) \texttt{FP32} accumulators. To eliminate the $3\times$ HBM/TMA traffic penalty of na\"ive $K$-stacking, we implement a 30-line patch to production \texttt{CUTLASS SM100 2SM} (\texttt{cta\_group::2}) \texttt{CollectiveMainloop} that holds paired pipeline stages $(s_0, s_1)$ co-resident in shared memory, reusing $A_0$ and $B_0$ in SMEM and $SF_{B0}$ in TMEM while cutting TMA traffic by $33\%$. Evaluated on live NVIDIA \texttt{GB200} silicon and all $48$ real layers of \texttt{Qwen3-30B-A3B}:
\begin{itemize}
    \item \textbf{Faster Than Native \texttt{BF16} End-to-End:} Patched \texttt{CUTLASS SM100 2SM} \texttt{Dual3\_SMEM} sustains \textbf{$2.28$--$2.62\text{ PF/s}$} ($1.18\times$--$1.54\times$ faster than \texttt{cuBLAS BF16}; $+14.7\%$--$+42.2\%$ faster than $K$-stacked \texttt{Dual3}), remaining \textbf{$1.01\times$--$1.46\times$ faster end-to-end} even including our vectorized \textbf{$5.58\text{ TB/s}$} native \texttt{cvt.e2m1x2} online quantizer. Asymmetric \texttt{Dual2} (\texttt{W4A(4+4)}) reaches \textbf{$3.44\text{ PF/s}$} ($1.80\times$--$1.92\times$ over \texttt{BF16}).
    \item \textbf{Live Hardware Real-Weight Accuracy Crushing \texttt{MXFP8}:} Executing real \texttt{Qwen3-30B-A3B} weights (\texttt{L00}--\texttt{L47}) live through patched \texttt{CUTLASS Dual3\_SMEM} silicon yields \textbf{$36.3\text{ dB}$} ($1\times$), \textbf{$38.0$--$38.6\text{ dB}$} ($50\times$), and \textbf{$39.9$--$40.3\text{ dB}$} ($200\times$) SQNR---beating live \texttt{CUTLASS MXFP8\_2SM} ($28.5\text{ dB}$) by \textbf{$+7.8$ to $+11.8\text{ dB}$} ($6.1\times$--$15.0\times$ lower MSE) and live \texttt{NVFP4\_1t} by \textbf{$+20.2\text{ dB}$} ($105\times$ lower MSE) with zero rotations.
    \item \textbf{All-48-Layer Frontier MoE \& 3-GEMM Training:} Across all 48 real \texttt{Qwen3-30B-A3B} layers, \texttt{Dual3} beats \texttt{MXFP8} on every metric (\textbf{KL $0.0309$ vs.\ $0.0636$, $2.06\times$ lower}; $12.17\times$ lower than \texttt{NVFP4+FWHT128}; Top-1 $90.5\%$ vs.\ $85.7\%$; MoE Top-8 Routing Jaccard $96.90\%$ vs.\ $94.65\%$) and outperforms \texttt{MXFP8} in 200-step all-3-GEMM ($Y, dX, dW$) outlier training.
\end{itemize}
\end{abstract}

\section{Introduction}
Blackwell's fifth-generation Tensor Cores (\texttt{tcgen05.mma}) introduce native hardware support for micro-scaled 4-bit floating-point (\texttt{NVFP4}: 1 sign, 2 exponent, 1 mantissa bit \texttt{E2M1} paired with an \texttt{E4M3} scale factor per $16$ contiguous elements). On \texttt{GB200} ($152$ SMs, \texttt{sm\_100a}), dense \texttt{NVFP4} sustains $9.65\text{ PF/s}$ bare-metal and up to $6.56$--$7.23\text{ PF/s}$ in production \texttt{CUTLASS 2SM} GEMMs---a $3.2\times$--$3.8\times$ empirical throughput advantage over \texttt{cuBLAS BF16} ($1.54$--$1.85\text{ PF/s}$ on realistic random matrices).

Yet \texttt{NVFP4} remains brittle for frontier LLM training, high-precision inference, and reasoning-heavy MoE models. Two structural obstacles explain why:
\begin{enumerate}
    \item \textbf{Intra-Block Outlier Masking:} A single $50\times$--$200\times$ channel spike forces the $1{\times}16$ block scale $s_0 = \max_i |x_i|/6$ upward by two orders of magnitude, quantizing the remaining $15$ elements to zero or $\pm 0.5 s_0$ ($\approx 17$--$20\text{ dB}$ SQNR).
    \item \textbf{The Hadamard Rotation Inversion on \texttt{NVFP4}:} While Walsh-Hadamard transforms (\texttt{FWHT}) rescue coarse $1{\times}32$ \texttt{UE8M0} \texttt{MXFP4}~\cite{tseng2024quip,ashkboos2024quarot}, applying \texttt{FWHT-128} to $1{\times}16$ \texttt{E4M3} \texttt{NVFP4} on real \texttt{Qwen3-30B-A3B} SwiGLU layers \emph{destroys} natural block sparsity and degrades 48-layer KL divergence by $2.15\times$ ($0.1747 \to 0.3760$, Top-1 $71.4\% \to 61.9\%$).
\end{enumerate}

We resolve both obstacles by observing that Blackwell's $4\times$ \texttt{NVFP4}/\texttt{BF16} Tensor-Core throughput ratio leaves enough compute headroom to execute \textbf{two or three micro-scaled \texttt{NVFP4} MMAs faster than a single \texttt{BF16} MMA}, provided the decomposition accumulates in-place inside Blackwell's $256\text{ KB/SM}$ Tensor Memory (\texttt{TMEM}) and reuses shared-memory operands across terms.

\section{Residual-Microscaled Dual-NVFP4 Theory}

\subsection{Two-Term Independent-Scale Decomposition}
Given a $16$-element block $\mathbf{x} \in \mathbb{R}^{16}$ and global tensor scale $g$, define the first-pass \texttt{NVFP4} quantization:
\begin{equation}
s_0 = Q_{\text{E4M3}}\!\left(\frac{\|\mathbf{x}\|_\infty}{6 g}\right) g, \quad \mathbf{q}_0 = Q_{\text{E2M1}}\!\left(\frac{\mathbf{x}}{s_0}\right), \quad \mathbf{a}_0 = s_0 \mathbf{q}_0.
\end{equation}
Define the residual $\mathbf{r} = \mathbf{x} - \mathbf{a}_0$ and quantize $\mathbf{r}$ with an \emph{independent} second $1{\times}16$ \texttt{E4M3} block scale $s_1$:
\begin{equation}
s_1 = Q_{\text{E4M3}}\!\left(\frac{\|\mathbf{r}\|_\infty}{6 g}\right) g, \quad \mathbf{q}_1 = Q_{\text{E2M1}}\!\left(\frac{\mathbf{r}}{s_1}\right), \quad \mathbf{a}_1 = s_1 \mathbf{q}_1.
\end{equation}

\begin{theorem}[Rotation-Free Outlier Immunity \& Effective Bit-Depth]
For any block $\mathbf{x} \in \mathbb{R}^{16}$ containing $k \ge 1$ outlier entries of magnitude $M \gg \sigma$ alongside bulk entries bounded by $\sigma$:
\begin{enumerate}
    \item Because \texttt{E2M1} spans $\{0, 0.5, 1, 1.5, 2, 3, 4, 6\}$ with relative step $\delta_{\max} \le \frac{1}{4}$ on $[1,6]$, the residual satisfies $\|\mathbf{r}\|_\infty \le \max\!\left(\frac{1}{6} M_{\text{rem}},\, \sigma\right)$ where the outlier's own quantization residual is bounded by $\le \frac{1}{6} M$ (and for isolated outliers aligned to $s_0 = M/6$, $|\mathbf{r}_{\text{out}}| \le \frac{1}{16} M$ due to \texttt{E4M3}'s 3-bit mantissa).
    \item When all elements are bounded uniformly ($\|\mathbf{x}\|_\infty = M$), $\|\mathbf{r}\|_\infty \le \frac{1}{6} M$, so $s_1 \le \frac{1}{6} s_0$, raising the worst-case quantization SNR by $\ge 20 \log_{10}(6) \approx +15.6\text{ dB}$ (empirical $+18.9$ to $+20.2\text{ dB}$, matching $7.5$--$8.0$ effective mantissa bits---exceeding \texttt{FP8 E4M3}'s $3$-bit mantissa!).
    \item For the 3-term truncated product $\widehat{AB} = A_0 B_0 + A_1 B_0 + A_0 B_1$, the dropped cross-term satisfies $\|A_1 B_1\|_F / \|A B\|_F \le \epsilon_A \epsilon_B \approx 2^{-7.5} \approx 5\times 10^{-3}$ ($\approx 38$--$40\text{ dB}$ SQNR ceiling), eliminating $25\%$ of MMAs with $<1.5\text{ dB}$ SQNR difference vs.\ the 4-term product.
\end{enumerate}
\end{theorem}

Crucially, because $A_0$ and $A_1$ share a single global tensor scale $g_A$ while \texttt{E4M3} spans $28{,}000\times$ dynamic range ($[2^{-6}, 448]$), $s_1$ stays cleanly in the normal \texttt{E4M3} range even when $s_1 / s_0 \approx 10^{-2}$.

\subsection{Why \texttt{Dual3} Beats \texttt{MXFP8} Despite Using 4-Bit Arithmetic}
An \texttt{MXFP8} block ($1{\times}32$, power-of-two \texttt{UE8M0} scale, $3$-bit \texttt{E4M3} mantissa) suffers up to $2\times$ wasted dynamic range from power-of-two scale rounding and a $3$-bit mantissa ($28.5\text{ dB}$ SQNR). \texttt{Dual-NVFP4} uses finer $1{\times}16$ blocks with $3$-bit mantissa \texttt{E4M3} scales \emph{twice}, yielding $\approx 36.3$--$40.3\text{ dB}$ SQNR ($+8$ to $+11.8\text{ dB}$ above \texttt{MXFP8}) at $9.0\text{ bits/elem}$ ($2 \times 4.5\text{ b}$), while executing on Blackwell's $9.65\text{ PF/s}$ \texttt{NVFP4} datapath.

\section{Blackwell \texttt{sm\_100a} System Architecture}

\subsection{5.58 TB/s Vectorized Online Dual-NVFP4 Quantizer}
Na\"ive warp-shuffle quantization with byte stores saturates at $0.30\text{ TB/s}$ ($355\,\mu\text{s}$ for $4\text{K}{\times}8\text{K}$). Our vectorized CUDA kernel (\texttt{k\_dual\_nvfp4\_quant\_vec}) assigns one thread per $16$-element block, issuing two coalesced $128$-bit \texttt{uint4} \texttt{BF16} loads, native Blackwell \texttt{\_\_nv\_cvt\_float2\_to\_fp4x2} and \texttt{\_\_nv\_cvt\_fp4x2\_to\_halfraw2} (\texttt{cvt.rn.f16x2.e2m1x2}) for zero-LUT residual computation, and coalesced $64$-bit \texttt{uint64\_t} stores for both $A_0$ and $A_1$. On \texttt{GB200}, this sustains \textbf{$5.19$--$5.58\text{ TB/s}$} ($20.2\,\mu\text{s}$ for $4\text{K}{\times}8\text{K}$, $39.8\,\mu\text{s}$ for $8\text{K}{\times}8\text{K}$, $150.3\,\mu\text{s}$ for $16\text{K}{\times}16\text{K}$)---a $18.6\times$ speedup over the shuffle baseline.

\subsection{SMEM-Reused \texttt{CUTLASS SM100 2SM} Mainloop Patch}
Computing $A_0 B_0 + A_1 B_0 + A_0 B_1$ via $3\times$ unfused GEMMs pays $2\times$ extra HBM read/write of $D$, while $3K$-concatenated single-launch GEMMs read $A_0$ and $B_0$ twice from HBM and TMA ($3\times$ TMA traffic).

Instead, we interleave $(A_0, A_1)$ and $(B_0, B_1)$ along $K$ in alternating $256$-element \texttt{CUTLASS} $K$-tiles (physical depth $2K$, $33\%$ less HBM/TMA traffic than $3K$) and apply a 30-line patch to \texttt{CUTLASS}'s \texttt{sm100\_blockscaled\_mma\_warpspecialized.hpp} (\texttt{KernelTmaWarpSpecialized2SmNvf4Sm100}, \texttt{Shape<\_256,\_256,\_256>}, \texttt{Cluster<\_2,\_1,\_1>}):
\begin{enumerate}
    \item Wait on pipeline stage $s_0 = (A_0, B_0, SF_{A0}, SF_{B0})$, peek stage $s_1$, copy $(SF_{A0}, SF_{B0})$ to \texttt{TMEM}, and immediately fire Term~1 ($A_0 B_0$) so stage $s_1$'s TMA load hides behind Term~1's $4$ \texttt{tcgen05} MMAs.
    \item Wait on stage $s_1 = (A_1, B_1, SF_{A1}, SF_{B1})$ while keeping $s_0$ pinned in SMEM; copy \emph{only} $SF_{A1}$ to \texttt{TMEM} and fire Term~2 ($A_1 B_0$), reusing both $B_0$ in SMEM and $SF_{B0}$ in \texttt{TMEM}.
    \item For \texttt{Dual3}, reload $(SF_{A0}, SF_{B1})$ and fire Term~3 ($A_0 B_1$), reusing $A_0$ in SMEM, accumulating all three terms in-place into the same \texttt{TMEM} \texttt{FP32} accumulators before releasing $(s_0, s_1)$.
\end{enumerate}
Critically, disabling CUTLASS's unrolled \texttt{IsOverlappingAccum} prologue avoids leaving an odd trailing $K$-tile that would deadlock the paired consumer loop.

\section{Live NVIDIA GB200 Silicon Evaluation}

\subsection{Bare-Metal Peak \& Production \texttt{CUTLASS 2SM} Roofline}
On all $152$ SMs ($76$ 2-SM clusters) of a live \texttt{GB200}, bare-metal SMEM-reused \texttt{tcgen05.mma.cta\_group::2} executes $1$-term \texttt{NVFP4} at $9.589\text{ PF/s}$ ($128.09\text{ cyc/tile}$), $2$-term \texttt{Dual2} at $4.825\text{ eff-PF/s}$ ($256.09\text{ cyc}$), and $3$-term \texttt{Dual3} at \textbf{$3.224\text{ eff-PF/s}$} ($384.09\text{ cyc}$).

Table~\ref{tab:cutlass} reports production CUDA-Graph-captured throughput on live \texttt{GB200} against realistic random \texttt{cuBLAS BF16}, verifying exact analytical outputs ($d_0 = 1.25 K$ for \texttt{Dual2}, $1.50 K$ for \texttt{Dual3} under $s_0{=}1.0, s_1{=}0.25$).

\begin{table}[t]
\centering
\caption{Production Patched \texttt{CUTLASS SM100 2SM} \texttt{Dual-NVFP4} Roofline on Live \texttt{GB200} (Effective Logical \texttt{TFLOPS}; \texttt{E2E} Includes Online \texttt{5.58 TB/s} Quantizer).}
\label{tab:cutlass}
\resizebox{\columnwidth}{!}{%
\begin{tabular}{lrrrrrrrr}
\toprule
\textbf{Shape} & \textbf{BF16} & \textbf{NV4\_1t} & \textbf{D2\_KStk} & \textbf{D3\_KStk} & \textbf{D3\_SMEM} & \textbf{D3\_E2E} & \textbf{SMEM/KStk} & \textbf{D3\_E2E/BF} \\
\midrule
$512{\times}8\text{K}{\times}8\text{K}$ & $1569.9$ & $4843.8$ & $2323.0$ & $1569.6$ & $\mathbf{1850.3}$ & $\mathbf{1582.2}$ & $+17.9\%$ & $\mathbf{1.01\times}$ \\
$4\text{K}{\times}8\text{K}{\times}8\text{K}$ & $1782.8$ & $6556.7$ & $3382.0$ & $2223.3$ & $\mathbf{2550.9}$ & $\mathbf{2326.1}$ & $+14.7\%$ & $\mathbf{1.31\times}$ \\
$4\text{K}{\times}14\text{K}{\times}4\text{K}$ & $1832.7$ & $5897.7$ & $3443.9$ & $1947.9$ & $\mathbf{2618.7}$ & $\mathbf{2452.7}$ & $+34.4\%$ & $\mathbf{1.34\times}$ \\
$4\text{K}^3$ & $1807.7$ & $5525.6$ & $2808.0$ & $1606.0$ & $\mathbf{2283.0}$ & $\mathbf{1892.1}$ & $+42.2\%$ & $\mathbf{1.05\times}$ \\
$8\text{K}^3$ & $1829.8$ & $5583.9$ & $2730.2$ & $1966.7$ & $\mathbf{2435.6}$ & $\mathbf{2213.8}$ & $+23.8\%$ & $\mathbf{1.21\times}$ \\
$16\text{K}^3$ & $1538.6$ & $5309.8$ & $2775.6$ & $1880.1$ & $\mathbf{2370.6}$ & $\mathbf{2241.4}$ & $+26.1\%$ & $\mathbf{1.46\times}$ \\
\bottomrule
\end{tabular}}
\end{table}

Our SMEM-reused \texttt{Dual3\_SMEM} beats $3K$ $K$-stacking by \textbf{$+14.7\%$ to $+42.2\%$} across all 6 shapes, surpasses \texttt{cuBLAS BF16} by \textbf{$1.18\times$--$1.54\times$} in GEMM ($2.28$--$2.62\text{ PF/s}$), and remains \textbf{$1.01\times$--$1.46\times$ faster end-to-end} including online quantization.

\subsection{Zero-Simulation Live Silicon Audit on Real \texttt{Qwen3-30B-A3B} Weights}
To eliminate any discrepancy between simulation and silicon, we quantized real \texttt{Qwen3-30B-A3B} weights (\texttt{L00, L15, L23, L31, L47} \texttt{q\_proj} and \texttt{o\_proj}) and $1\times$/$50\times$/$200\times$ outlier activations on GPU via \texttt{cvt.e2m1x2} directly into CUTLASS's MMA scale layouts and executed live hardware \texttt{CUTLASS SM100 2SM} kernels against \texttt{FP32} ground truth (Table~\ref{tab:live_hw}).

\begin{table}[t]
\centering
\caption{Live Hardware \texttt{CUTLASS SM100 2SM} SQNR (dB) on Real \texttt{Qwen3-30B-A3B} Weights ($M{=}4096$, Zero Simulation).}
\label{tab:live_hw}
\resizebox{\columnwidth}{!}{%
\begin{tabular}{llrrrrrr}
\toprule
\textbf{Real Tensor} & \textbf{Outlier} & \textbf{BF16} & \textbf{MXFP8} & \textbf{NV4\_1t} & \textbf{Dual2} & \textbf{Dual3} & \textbf{D3 vs FP8} \\
\midrule
\texttt{L00\_o\_proj} & $1\times$ & $52.59$ & $28.51$ & $17.48$ & $20.47$ & $\mathbf{36.34}$ & $\mathbf{+7.83\text{ dB}}$ \\
\texttt{L00\_o\_proj} & $50\times$ & $52.61$ & $28.50$ & $19.11$ & $20.67$ & $\mathbf{38.64}$ & $\mathbf{+10.14\text{ dB}}$ \\
\texttt{L00\_o\_proj} & $200\times$ & $52.61$ & $28.54$ & $20.10$ & $20.63$ & $\mathbf{40.30}$ & $\mathbf{+11.76\text{ dB}}$ \\
\texttt{L23\_o\_proj} & $50\times$ & $52.61$ & $28.52$ & $18.70$ & $20.37$ & $\mathbf{38.01}$ & $\mathbf{+9.50\text{ dB}}$ \\
\texttt{L31\_q\_proj} & $200\times$ & $52.60$ & $28.48$ & $19.79$ & $20.34$ & $\mathbf{39.90}$ & $\mathbf{+11.41\text{ dB}}$ \\
\texttt{L47\_o\_proj} & $200\times$ & $52.62$ & $28.50$ & $19.88$ & $20.44$ & $\mathbf{40.13}$ & $\mathbf{+11.64\text{ dB}}$ \\
\bottomrule
\end{tabular}}
\end{table}

Across all $15/15$ real tensor$\times$outlier configurations, live \texttt{Dual3\_SMEM} silicon beats live \texttt{CUTLASS MXFP8\_2SM} by \textbf{$+7.81$ to $+11.76\text{ dB}$} ($6.1\times$--$15.0\times$ lower MSE) and live \texttt{NVFP4\_1t} by up to \textbf{$+20.2\text{ dB}$} ($105\times$ lower MSE), sustaining $1{,}782$--$1{,}851\text{ TF/s}$ even on small $[4096,2048,4096]$ shapes.

\subsection{All-48-Layer \texttt{Qwen3-30B-A3B} Audit \& 200-Step Outlier Training}
Evaluating all $48$ real layers of \texttt{Qwen3-30B-A3B} (attention + 128-expert MoE + \texttt{lm\_head}, Table~\ref{tab:48l}) confirms that \texttt{Dual3} dominates every quantized baseline without rotations: $2.06\times$ lower KL than \texttt{MXFP8}, $3.18\times$ lower than \texttt{MXFP6-E2M3}, $5.65\times$ lower than \texttt{NVFP4\_1t}, and $12.17\times$ lower than \texttt{NVFP4+FWHT128}. Adaptively skipping $57.6\%$ of activation and $48.5\%$ of weight low-residual $1{\times}16$ blocks (\texttt{Dual3\_Adapt50}, $\approx 1.94$ effective FP4 MMAs) still beats \texttt{MXFP6} ($0.0767$ vs.\ $0.0982$ KL).

\begin{table}[t]
\centering
\caption{All-48-Layer Real \texttt{Qwen3-30B-A3B} Forward Pass Audit vs.\ \texttt{BF16}.}
\label{tab:48l}
\resizebox{\columnwidth}{!}{%
\begin{tabular}{lrrrr}
\toprule
\textbf{Quantization Mode} & \textbf{48L KL} & \textbf{Top-1} & \textbf{Logit SQNR} & \textbf{MoE Top-8 Jaccard} \\
\midrule
\textbf{\texttt{Dual3} (0 Hadamard)} & $\mathbf{0.030903}$ & $\mathbf{90.5\%}$ & $\mathbf{25.60\text{ dB}}$ & $\mathbf{96.90\%}$ \\
\texttt{MXFP8} ($1{\times}32$) & $0.063627$ & $85.7\%$ & $23.85\text{ dB}$ & $94.65\%$ \\
\texttt{Dual3\_Adapt50} & $0.076669$ & $81.0\%$ & $20.78\text{ dB}$ & $93.06\%$ \\
\texttt{MXFP6-E2M3} ($1{\times}32$) & $0.098220$ & $90.5\%$ & $18.07\text{ dB}$ & $94.21\%$ \\
\texttt{Dual2\_W4A44} ($4.5\text{ b/w}$) & $0.145552$ & $81.0\%$ & $17.41\text{ dB}$ & $90.66\%$ \\
\texttt{NVFP4\_1t} ($1{\times}16$) & $0.174679$ & $71.4\%$ & $15.65\text{ dB}$ & $88.24\%$ \\
\texttt{NVFP4\_FWHT128} & $0.376016$ & $61.9\%$ & $11.93\text{ dB}$ & $86.59\%$ \\
\bottomrule
\end{tabular}}
\end{table}

Finally, in 200-step all-3-GEMM ($Y, dX, dW$) training on real \texttt{Qwen3} expert weights with stochastic rounding on residual $A_1$, \texttt{Dual3\_SR} beats \texttt{MXFP8} at every outlier regime: final loss $0.00044$ vs.\ $0.00134$ ($1\times$), $0.00096$ vs.\ $0.00115$ ($50\times$), and $0.00338$ vs.\ $0.00348$ ($200\times$), reducing loss by $29\times$--$74\times$ over single-term \texttt{NVFP4\_SR} ($0.0218$--$0.0321$).

\section{Related Work \& Limitations}
Prior FP4 quantization methods rely on offline/online Hadamard rotations (\texttt{QuaRot}~\cite{ashkboos2024quarot}, \texttt{QuIP\#}~\cite{tseng2024quip}, \texttt{SpinQuant}), low-rank FP16 outlier branches (\texttt{SVDQuant}), or \texttt{FP8}/\texttt{MXFP6} fallbacks. \texttt{Dual-NVFP4} is orthogonal and strictly hardware-native: every term executes on Blackwell's $9.65\text{ PF/s}$ \texttt{NVFP4} \texttt{tcgen05} pipe with zero rotations and zero FP16 side-branches. \emph{Limitations:} \texttt{Dual3} uses $9.0\text{ b/elem}$ storage (vs.\ $4.5\text{ b}$ for 1-term \texttt{NVFP4}, matching \texttt{FP8} storage), making it optimal for compute-bound prefill, training, and outlier-sensitive layers; memory-bound $B{=}1$ decode should pair $4.5\text{ b/w}$ weights with \texttt{Dual2\_W4A44} or surgical \texttt{Dual3} on attention/router layers.

\section{Conclusion}
\texttt{Dual-NVFP4} establishes a new Pareto frontier on NVIDIA Blackwell: a 30-line SMEM-reused \texttt{CUTLASS SM100 2SM} patch that outperforms native \texttt{cuBLAS BF16} by $1.21\times$--$1.46\times$ end-to-end while outperforming \texttt{MXFP8} by $+7.8$ to $+11.8\text{ dB}$ SQNR and $2.06\times$ lower 48-layer KL divergence with zero Hadamard rotations.

\begin{thebibliography}{9}
\bibitem{ashkboos2024quarot} S.~Ashkboos et al., ``QuaRot: Outlier-Free 4-Bit Inference in Rotated LLMs,'' \emph{NeurIPS}, 2024.
\bibitem{tseng2024quip} A.~Tseng et al., ``QuIP\#: Even Better LLM Quantization with Hadamard Incoherence and Lattice Codebooks,'' \emph{ICML}, 2024.
\bibitem{nvidia2025blackwell} NVIDIA Corporation, ``Blackwell Architecture \& PTX ISA 8.7 (\texttt{tcgen05}),'' 2025.
\bibitem{cutlass2025} NVIDIA CUTLASS Team, ``CUTLASS 3.x SM100 Block-Scaled Warp-Specialized GEMM,'' 2025.
\end{thebibliography}
\end{document}
